\chapter{Results}\label{chapter:results}

\section{Morphological Inflection}

This section reports test-set performance on the morphological inflection task using the UniMorph-based datasets described in Chapter~\ref{chapter:datasets}. Each column corresponds to one language, identified by its ISO code (see Table~\ref{tab:languages-um}). Results are summarized in Table~\ref{tab:results-inflection-matrix}, which reports both accuracy and mean Levenshtein distance.
\begin{table}[h]
\centering
\normalsize
\setlength{\tabcolsep}{5pt}
\renewcommand{\arraystretch}{1.10}
\begin{tabular}{lrrrrrrrrr}
\toprule
Model & amh & azg & dsb & eng & grc & ita & kat & por & avg \\
\midrule
\multicolumn{10}{l}{\textit{Accuracy (\%, $\uparrow$)}} \\
Non-neural & 39.0 & 6.1 & \textbf{85.2} & 95.8 & 29.0 & 76.2 & 84.2 & 96.6 & 64.0 \\
Neural & \textbf{88.7} & 20.1 & 79.1 & 96.6 & 19.4 & 78.3 & \textbf{89.2} & \textbf{97.6} & 71.1 \\
ByT5 & 82.2 & \textbf{26.5} & 83.9 & 97.1 & \textbf{41.6} & \textbf{95.8} & 85.1 & 95.0 & \textbf{75.9} \\
ByT5 bi & 60.8 & 24.4 & 78.3 & \textbf{97.6} & 34.8 & 89.8 & 82.9 & 91.1 & 70.0 \\
\midrule
\multicolumn{10}{l}{\textit{Mean Levenshtein distance ($\downarrow$)}} \\
Non-neural & 1.010 & 5.660 & \textbf{0.267} & 0.086 & 1.852 & 0.761 & 0.430 & 0.054 & 1.265 \\
Neural & \textbf{0.162} & 3.703 & 0.379 & 0.069 & 2.212 & 0.266 & \textbf{0.297} & \textbf{0.031} & 0.890 \\
ByT5 & 0.237 & \textbf{3.413} & 0.287 & 0.067 & \textbf{1.713} & \textbf{0.091} & 0.480 & 0.072 & \textbf{0.795} \\
ByT5 bi & 0.606 & 3.564 & 0.375 & \textbf{0.050} & 2.020 & 0.243 & 0.634 & 0.149 & 0.955 \\
\bottomrule
\end{tabular}
\caption{Inflection results (test set): accuracy and mean Levenshtein distance.}
\label{tab:results-inflection-matrix}
\end{table}

Accuracy measures exact-match correctness of the predicted inflected form, while the mean Levenshtein distance captures the average character-level edit distance between predictions and gold forms (lower is better). Throughout this chapter, boldface indicates the best model for each language within each metric; in all languages in Table~\ref{tab:results-inflection-matrix}, the same model is best under both accuracy and mean Levenshtein distance.

Overall, the best-performing models are split between the neural baseline and the ByT5 variants, suggesting that both strong task-specific architectures and large pretrained sequence-to-sequence models can be competitive depending on the language. Moreover, the bidirectional ByT5 variant is trained on a 50/50 mixture of forward and inverse UniMorph examples (Section~\ref{sec:bidirectional}), yet it remains competitive in the forward direction and even outperforms the neural baseline for several languages (e.g., \texttt{azg}, \texttt{eng}, \texttt{grc}, \texttt{ita}). A notable exception is Lower Sorbian (\texttt{dsb}), where the non-neural baseline achieves the best accuracy and the lowest mean Levenshtein distance, indicating that well-designed rule-based approaches can still match or outperform neural systems in some settings.

\section{Morphological Analysis}

This section reports test-set performance on morphological analysis. Each column corresponds to one language, identified by its ISO code (see Table~\ref{tab:languages-um}). As in the evaluation pipeline, results are reported for (i) predicting the lemma and (ii) predicting the target \ac{MSD} (interpreted as a set of tags). Since the inverse and bidirectional ByT5 models are trained and tested on UniMorph, while the context-based and LLM experiments use UD-derived data, results are reported separately for the two settings.

UniMorph-based results for the ByT5 inverse models are shown in Table~\ref{tab:results-analysis-um-lemma-matrix} (lemma metrics) and Table~\ref{tab:results-analysis-um-msd-matrix} (\ac{MSD} metrics).
\begin{table}[h]
\centering
\small
\setlength{\tabcolsep}{5pt}
\renewcommand{\arraystretch}{1.10}
\begin{tabular}{lrrrrrrrr}
\toprule
Model & amh & azg & dsb & eng & grc & ita & kat & por \\
\midrule
\multicolumn{9}{l}{\textit{Lemma accuracy (\%, $\uparrow$)}} \\
ByT5 inv & \textbf{76.2} & 34.4 & 61.3 & \textbf{97.0} & 59.7 & \textbf{94.5} & 75.6 & \textbf{94.8} \\
ByT5 bi-inv & 64.6 & \textbf{37.1} & \textbf{64.3} & 96.8 & \textbf{64.0} & 92.6 & \textbf{76.2} & 94.2 \\
\midrule
\multicolumn{9}{l}{\textit{Lemma mean Levenshtein ($\downarrow$)}} \\
ByT5 inv & \textbf{0.287} & 2.618 & 0.597 & \textbf{0.044} & 0.737 & \textbf{0.099} & 0.462 & \textbf{0.076} \\
ByT5 bi-inv & 0.408 & \textbf{2.343} & \textbf{0.546} & 0.048 & \textbf{0.701} & 0.141 & \textbf{0.434} & 0.091 \\
\bottomrule
\end{tabular}
\caption{Morphological analysis (UniMorph, test set): lemma metrics.}
\label{tab:results-analysis-um-lemma-matrix}
\end{table}

\begin{table}[h]
\centering
\small
\setlength{\tabcolsep}{5pt}
\renewcommand{\arraystretch}{1.10}
\begin{tabular}{lrrrrrrrr}
\toprule
Model & amh & azg & dsb & eng & grc & ita & kat & por \\
\midrule
\multicolumn{9}{l}{\textit{\acs{MSD} accuracy (\%, $\uparrow$)}} \\
ByT5 inv & \textbf{83.2} & \textbf{77.8} & \textbf{46.2} & 81.3 & \textbf{73.4} & \textbf{78.9} & 72.7 & \textbf{66.1} \\
ByT5 bi-inv & 72.3 & 69.2 & 36.8 & \textbf{81.7} & 63.8 & 78.1 & \textbf{73.0} & 65.7 \\
\midrule
\multicolumn{9}{l}{\textit{\acs{MSD} F1 (\%, $\uparrow$)}} \\
ByT5 inv & \textbf{92.8} & \textbf{94.4} & \textbf{72.4} & 96.4 & \textbf{90.6} & \textbf{93.0} & \textbf{87.0} & \textbf{89.0} \\
ByT5 bi-inv & 88.4 & 89.0 & 67.9 & \textbf{96.6} & 84.9 & 92.5 & 84.9 & 88.5 \\
\bottomrule
\end{tabular}
\caption{Morphological analysis (UniMorph, test set): \acs{MSD} metrics.}
\label{tab:results-analysis-um-msd-matrix}
\end{table}


For UniMorph-based analysis, the model that achieves the best lemma accuracy for a given language is also best under lemma mean Levenshtein distance (Table~\ref{tab:results-analysis-um-lemma-matrix}). Notably, the bidirectional inverse model (ByT5 bi-inv), which is trained on a 50/50 mix of forward and inverse examples (Section~\ref{sec:bidirectional}), achieves the best lemma results in exactly half of the languages (\texttt{azg}, \texttt{dsb}, \texttt{grc}, \texttt{kat}), indicating that splitting supervision across the two directions does not necessarily harm lemma prediction. For \ac{MSD} prediction, ByT5 inv yields the strongest results in most languages; the winner is consistent between \ac{MSD} accuracy and F1 in all but one case (\texttt{kat}), where the two metrics select different models (Table~\ref{tab:results-analysis-um-msd-matrix}).

UD-based results for the context-based ByT5 model and the LLM baselines are shown in Table~\ref{tab:results-analysis-ud-lemma-matrix} (lemma metrics) and Table~\ref{tab:results-analysis-ud-msd-matrix} (\ac{MSD} metrics).
\begin{table}[h]
\centering
\small
\setlength{\tabcolsep}{5pt}
\renewcommand{\arraystretch}{1.10}
\begin{tabular}{lrrrrrr}
\toprule
Model & amh & eng & grc & ita & kat & por \\
\midrule
\multicolumn{7}{l}{\textit{Lemma accuracy (\%, $\uparrow$)}} \\
ByT5 ctx & \textbf{89.9} & \textbf{96.9} & 14.4 & \textbf{95.5} & \textbf{6.1} & \textbf{74.5} \\
LLM Llama & 66.7 & 66.4 & \textbf{35.1} & 78.2 & 4.9 & 71.3 \\
LLM Qwen & 47.5 & 49.1 & 32.7 & 42.8 & 4.0 & 51.6 \\
\midrule
\multicolumn{7}{l}{\textit{Lemma mean Levenshtein ($\downarrow$)}} \\
ByT5 ctx & \textbf{0.131} & \textbf{0.072} & \textbf{4.030} & \textbf{0.149} & \textbf{3.188} & \textbf{0.596} \\
LLM Llama & 0.798 & 0.855 & 4.057 & 2.002 & 4.750 & 0.751 \\
LLM Qwen & 4.323 & 5.465 & 6.077 & 8.332 & 17.819 & 8.453 \\
\bottomrule
\end{tabular}
\caption{Morphological analysis (UD, test set): lemma metrics.}
\label{tab:results-analysis-ud-lemma-matrix}
\end{table}


As discussed in Section~\ref{sec:analysis-eval}, for the UD-based setting \ac{MSD} metrics are only reported for languages whose UD-derived labels could be converted into the evaluation tag inventory; therefore, Amharic (\texttt{amh}) and Georgian (\texttt{kat}) are omitted from Table~\ref{tab:results-analysis-ud-msd-matrix}.
\begin{table}[h]
\centering
\small
\setlength{\tabcolsep}{5pt}
\renewcommand{\arraystretch}{1.10}
\begin{tabular}{lrrrrr}
\toprule
Model & eng & grc & ita & por & avg \\
\midrule
\multicolumn{6}{l}{\textit{\acs{MSD} accuracy (\%, $\uparrow$)}} \\
ByT5 ctx & \textbf{95.7} & 13.6 & \textbf{96.8} & \textbf{71.8} & \textbf{69.5} \\
LLM Llama & 10.5 & 36.6 & 3.5 & 1.9 & 13.1 \\
LLM Qwen & 6.7 & \textbf{47.8} & 21.8 & 2.8 & 19.8 \\
\midrule
\multicolumn{6}{l}{\textit{\acs{MSD} F1 (\%, $\uparrow$)}} \\
ByT5 ctx & \textbf{98.3} & 30.6 & \textbf{98.4} & \textbf{83.8} & \textbf{77.8} \\
LLM Llama & 37.5 & 81.7 & 17.7 & 45.8 & 45.7 \\
LLM Qwen & 40.7 & \textbf{81.8} & 24.7 & 41.5 & 47.2 \\
\bottomrule
\end{tabular}
\caption{Morphological analysis (UD, test set): \acs{MSD} metrics.}
\label{tab:results-analysis-ud-msd-matrix}
\end{table}


In the UD-based setting, the context-aware ByT5 model achieves the strongest performance for most languages and metrics (Tables~\ref{tab:results-analysis-ud-lemma-matrix} and \ref{tab:results-analysis-ud-msd-matrix}), substantially outperforming the prompted LLM baselines. Ancient Greek is a notable exception: for \texttt{grc}, the LLM baselines outperform the context-aware model on lemma accuracy and, even more clearly, on \ac{MSD} accuracy and F1. At the same time, the lemma mean Levenshtein distances for \texttt{grc} remain high across all models, suggesting that the UD-derived Greek setting is particularly challenging in this configuration.

Two dataset-specific effects are worth noting when interpreting language-level outliers. First, in the prepared UniMorph subset used in this thesis, the Ancient Greek (\texttt{grc}) splits contain nominal paradigms only (nouns and adjectives; no verb instances), whereas the UD-derived analysis experiments are verb-focused; consequently, Greek results are not directly comparable across the two data sources. This mismatch is a plausible contributing factor to why \texttt{grc} behaves differently in the UD-based tables, including the observation that the LLM baselines outperform the context-aware ByT5 model for Greek, despite the strong UniMorph-based results for ByT5 on \texttt{grc}. Second, Georgian exhibits a systematic mismatch in lemma conventions between UniMorph and UD: in the UniMorph Georgian data, verbs are grouped under a UniMorph-specific citation form (see also \parencite{guriel-etal-2022-morphological}), while UD lemmas for Georgian are typically based on the \textit{masdar} (verbal noun) form.\footnote{UD lemmatization convention is described in the UD\_Georgian-GNC README: \url{https://github.com/UniversalDependencies/UD_Georgian-GNC/blob/master/README.md}.} This difference can substantially affect lemma-level scores and makes cross-source comparisons for Georgian particularly sensitive to annotation conventions.
